\begin{textsample}{NEG Dim 2 – 2021\_c5 – Score -2.00 – 2021\_c5/enade\_2021\_letras\_por\_d96b21bbf0.txt}  \label{ex:f2_neg_145}
CAPÍTULO 5 DISTRIBUIÇÃO DOS CONCEITOS 5.1 PANORAMA NACIONAL DA DISTRIBUIÇÃO DOS CONCEITOS O Conceito Enade é um indicador de qualidade que avalia os cursos por intermédio dos desempenhos dos estudantes a partir dos resultados do Enade.
Os valores possíveis vão de 1 ( pior situação ) a 5 ( melhor situação ).
Os cursos com apenas um ou sem nenhum concluinte participante não obtêm o Conceito Enade, ficando Sem Conceito ( SC ).
Na Tabela 5.1, são apresentadas a quantidade e a distribuição de cursos de Letras-Português ( Licenciatura ) participantes do Enade/2021, por faixa de conceito e Grande Região. \textbf{Enfatiza} -se, mais uma vez, que a diferença entre os cursos tabulados neste capítulo e no Capítulo 2 corresponde aos cursos que não foram avaliados, ou seja, àqueles, em princípio, sem alunos concluintes que tivessem sido inscritos para a prova.
Observando -se os dados da Tabela 5.1, nota -se que, dos 292 cursos participantes, 104 ( 35,6 \% ) classificaram -se com conceito 3, o valor modal.
Esse foi também o conceito modal na maioria das regiões : Nordeste ( 40,0 \% ), Sudeste ( 28,6 \% ), Centro-Oeste ( 42,3 \% ) e Sul ( 30,0 \% ).
Nessa última, a distribuição foi bimodal, com o conceito 4 na mesma proporção.
O conceito 2 foi o segundo mais frequente em nível nacional ( 28,8 \% ), correspondendo a 84 cursos, sendo o conceito modal na região Norte, com 12 cursos ( 44,4 \% ).
Já o conceito 4 foi o terceiro mais frequente ( 17,8 \%, correspondendo a 52 cursos ).
Houve, ainda, 26 cursos ( 8,9 \% ) que receberam o conceito 5, e 17 cursos ( 5,8 \% ) que receberam o conceito 1.
Dos 292 cursos de Letras-Português ( Licenciatura ), nove ( 3,1 \% ) ficaram Sem Conceito ( SC ).
A região Norte participou com 27 cursos, o equivalente a 9,2 \% do total nacional.
Desses, 12 cursos receberam o conceito 2, o que equivale a 44,4 \% do total regional, o conceito modal na região, como já comentado.
A dois cursos ( 7,4 \% ) atribuiu -se o conceito 1, a 11 cursos ( 40,7 \% ), o conceito 3, e dois cursos ( 7,4 \% ) receberam o conceito 4.
Nenhum curso ficou Sem Conceito ( SC ) ou recebeu o conceito 5 nessa região.
A região Nordeste participou com 115 cursos, o equivalente a 39,4 \% do total nacional.
Nessa região, 46 cursos ( 40,0 \% em termos regionais ) obtiveram o conceito 3, o conceito modal para a região, como já comentado.
Os conceitos 1 e 2 foram atribuídos, respectivamente, a 13 e a 39 cursos ( 11,3 \% e 33,9 \%, respectivamente ).
Ainda nessa região, 11 cursos ( 9,6 \% ) foram avaliados com o conceito 4, e quatro cursos ( 3,5 \% ) foram avaliados com o conceito 5.
Dois cursos ficaram Sem Conceito ( SC ).
Dos 84 cursos participantes da região Sudeste ( 28,8 \% do total nacional ), o conceito 3 foi o conceito modal, como já anteriormente assinalado, tendo sido obtido por 24 cursos ( 28,6 \% ).
O conceito 2 foi atribuído a 18 cursos ( 21,4 \% ), e o conceito 4, a 23 cursos ( 27,4 \% ).
O conceito 5 foi atribuído a 16 cursos ( 19,0 \% ).
Nessa região, três cursos ficaram Sem Conceito ( SC ), e nenhum recebeu o conceito 1.
Os 40 cursos da região Sul corresponderam a 13,7 \% do total nacional.
Como já indicado, a distribuição, nessa região, foi bimodal, com os conceitos 3 e 4 concentrando 12 cursos cada, e correspondendo, juntos, a 60,0 \% do total regional.
Os conceitos 1 e 2 foram atribuídos a um ( 2,5 \% ) e a oito cursos ( 20,0 \% ), respectivamente, e o conceito 5, a cinco cursos ( 12,5 \% ).
Dois cursos ficaram Sem Conceito ( SC ) na região Sul.
Já dos 26 cursos participantes na região Centro-Oeste ( 8,9 \% do total nacional ), 11 ( 42,3 \% em termos regionais ) receberam conceito 3, o conceito modal, como já destacado.
Um curso ( 3,8 \% ) obteve o conceito 1, sete cursos ( 26,9 \% ) obtiveram o conceito 2, quatro cursos ( 15,4 \% ) obtiveram o conceito 4 e um ( 3,8 \% ) recebeu o conceito 5.
Nessa região, dois cursos ficaram Sem Conceito ( SC ).
No Gráfico 5.1, é apresentada a distribuição cumulativa do Conceito Enade, segundo a Grande Região.
Em um gráfico de distribuição cumulativa, em cada valor de abscissa ( eixo x ), o valor apresentado no eixo das ordenadas ( eixo y ) é o valor acumulado da variável ( Conceito Enade ).
Observando -se a linha vertical pontilhada, que ajuda a visualizar o percentual de conceitos acumulados até o valor 3, ou seja, os conceitos 1, 2 e 3, pode -se notar que, por exemplo, a região Norte ( linha azul ) apresentou 92,6 \% dos cursos com conceito 3 ou abaixo desse valor.
A Sudeste ( linha vermelha ), por outro lado, apresentou 51,9 \% dos cursos na mesma situação : com conceito 3 ou abaixo.
Destaca -se que em algumas regiões o conceito 5 pode não ter sido alcançado por nenhum curso.
Foi o que ocorreu na região Norte ( linha azul ), na qual todos os cursos obtiveram conceitos menores ou iguais a 4, com a poligonal alcançando 100 \% no conceito 4.
Em linhas gerais, em um gráfico de distribuição cumulativa, poligonais mais à esquerda representam as Grandes Regiões com a pior distribuição, e poligonais mais à direita, as Grandes Regiões com os melhores conceitos.
A região Norte ( linha azul ) e a Nordeste ( linha verde ) apresentaram os cursos com os conceitos mais baixos, ao passo que as regiões Sul ( linha rosa ) e Sudeste ( linha vermelha ) apresentaram a distribuição com os valores maiores.
Cursos em IES situadas na região Centro-Oeste ( linha laranja ) apresentaram uma situação intermediária. 5.2 CONCEITOS POR CATEGORIA ADMINISTRATIVA, POR MODALIDADE DE ENSINO E GRANDE REGIÃO Na Tabela 5.2, é apresentada a distribuição dos cursos de Letras-Português ( Licenciatura ) participantes do Enade/2021, por Categoria Administrativa e Modalidade de Ensino, de acordo com os conceitos por eles alcançados, segundo as Grandes Regiões brasileiras.
Dos 292 cursos participantes, 193 ( 66,1 \% ) eram ministrados em Instituições Públicas, e 99 ( 33,9 \% ), em Instituições Privadas.
Quanto à Modalidade de Ensino, 68 cursos eram a Distância, e 224, Presenciais.
De acordo com as informações da Tabela 5.2, em termos nacionais, entre Instituições Públicas e Privadas, observa -se que, dos 26 cursos avaliados com conceito 5, 16 eram oferecidos em IES Públicas, e 10, em IES Privadas.
Dos 193 cursos participantes de IES Públicas, o conceito 3 foi o valor modal, atribuído a 71 cursos ( 36,8 \% ).
Dos demais cursos de Instituições Públicas participantes, 13 cursos ( 6,7 \% ) foram avaliados com o conceito 1, 54 cursos ( 28,0 \% ), com o conceito 2, 38 cursos ( 19,7 \% ), com o conceito 4 e, como já comentado, 16 cursos ( 8,3 \% ), com o conceito 5.
Nessa categoria, um curso ficou Sem Conceito ( SC ).
Na rede privada, o conceito modal foi o 3, atribuído a 33 cursos, o correspondente a 33,3 \% dos 99 cursos da categoria.
Dos demais cursos participantes, quatro ( 4,0 \% ) receberam o conceito 1, e 30 ( 30,3 \% ), o conceito 2.
O conceito 4 foi atribuído a 14 cursos ( 14,1 \% ), e o conceito 5, a dez cursos ( 10,1 \% ).
Nessa Categoria Administrativa, oito cursos ( 8,1 \% ) ficaram Sem Conceito ( SC ).
Dos 224 cursos participantes da modalidade presencial, o conceito 3 foi o valor modal, atribuído a 82 cursos ( 36,6 \% ).
Dos demais cursos nessa Modalidade de Ensino, 12 cursos ( 5,4 \% ) receberam o conceito 1, 58 cursos ( 25,9 \% ) receberam o conceito 2, 42 cursos ( 18,8 \% ), o conceito 4, e 23 cursos ( 10,3 \% ), o conceito 5.
Nessa Modalidade de Ensino, sete cursos ( 3,1 \% ) ficaram Sem Conceito ( SC ).
Foram 68 os cursos a Distância participantes, dos quais 26 ( 38,2 \% ) obtiveram o conceito modal 2.
Nessa Modalidade de Ensino, cinco cursos ( 7,4 \% ) receberam o conceito 1 ; 22 cursos ( 32,4 \% ), o conceito 3 ; dez cursos ( 14,7 \% ), o conceito 4 ; e três cursos ( 4,4 \% ), o conceito 5.
Nessa região, dois cursos ( 2,9 \% ) ficaram Sem Conceito ( SC ).
Na análise por região, observa -se que, na região Norte, as Instituições Privadas participaram com quatro cursos ( 14,8 \% do total regional ).
O conceito modal para as Instituições Privadas na região foi o 2, atribuído a três cursos, correspondendo a 75,0 \% dessa categoria na região.
O curso restante foi avaliado com o conceito 3.
Nenhum ficou Sem Conceito ( SC ) ou recebeu os demais conceitos, nessa categoria e região.
As Instituições Públicas participaram com 23 cursos na região Norte ( 85,2 \% do total regional ), e o conceito modal foi o 3, atribuído a dez cursos.
Dos restantes, dois cursos receberam o conceito 1, nove cursos, o conceito 2, e dois cursos, o conceito 4.
Nenhum curso vinculado a Instituições Públicas ficou Sem Conceito ( SC ) ou recebeu conceito 5 na região Norte.
Todos os cursos oferecidos na região Norte foram Presenciais.
Na região Nordeste, dos 115 cursos participantes, 17 eram vinculados à rede privada, o equivalente a 14,8 \% do total da região.
Nessa região, a distribuição foi bimodal, com os conceitos 2 e 3 atribuídos a cinco cursos cada ( correspondendo, esses dez cursos, a 58,8 \% do total regional ).
Ainda nessa categoria e região, três cursos foram avaliados com o conceito 1, dois cursos, com o conceito 4, sendo que dois cursos ficaram Sem Conceito ( SC ) e nenhum recebeu o conceito 5.
As Instituições Públicas da região Nordeste contaram com 98 cursos ( 85,2 \% do total da região ).
Desses, 41 obtiveram conceito 3, o valor modal.
Os demais foram avaliados com os seguintes conceitos : 1 ( dez cursos ), 2 ( 34 cursos ), 4 ( nove cursos ) e 5 ( quatro cursos ).
Nenhum curso ficou Sem Conceito ( SC ).
Dos 95 cursos ( 82,6 \% ) participantes da Modalidade Presencial na região Nordeste, o conceito 3 foi o valor modal, atribuído a 39 cursos ( 41,1 \% ).
Dos demais cursos nessa Modalidade de Ensino, nove cursos ( 9,5 \% ) receberam o conceito 1, e 31 cursos ( 32,6 \% ), o conceito 2.
O conceito 4 foi atribuído a dez cursos ( 10,5 \% ), e o conceito 5, a quatro cursos ( 4,2 \% ).
Nessa Modalidade de Ensino, dois cursos ficaram Sem Conceito ( SC ).
Dos 20 cursos a Distância ( 17,4 \% ) na região Nordeste, oito cursos ( 40,0 \% ) receberam o conceito modal 2, quatro cursos ( 20,0 \% ), o conceito 1, sete cursos ( 35,0 \% ), o conceito 3, e um curso ( 5,0 \% ) recebeu o conceito 4.
Nenhum curso nessa Modalidade de Ensino na região Nordeste ficou Sem Conceito ( SC ) ou recebeu o conceito 5.
Na região Sudeste, a proporção de cursos da rede privada, 58,3 \%, correspondendo a 49 dos 84 cursos participantes, foi mais \textbf{elevada} do que a registrada nas demais regiões brasileiras.
Nessa categoria e região, a distribuição foi bimodal, com os conceitos 2 e 3 concentrados em 15 cursos cada ( correspondendo, esses 30 cursos, a 61,2 \% total regional ).
Os demais foram avaliados com o conceito 4 ( nove cursos ) e o conceito 5 ( sete cursos ).
Nessa combinação de categoria e região, três cursos ficaram Sem Conceito ( SC ), e nenhum recebeu o conceito 1.
Dos 35 cursos oferecidos por Instituições Públicas na região Sudeste, o conceito modal foi o 4, atribuído a 14 cursos.
Nessa região, três cursos receberam o conceito 2, nove cursos, o conceito três e nove cursos, o conceito 5, não havendo cursos Sem Conceito ( SC ) ou com conceito 1.
Também na região Sudeste, a maior parte dos cursos foi oferecida na modalidade Presencial : 55 de 84 cursos.
Para esses, a distribuição foi bimodal, com os conceitos 3 e 4 sendo atribuídos a 15 cursos cada.
Os demais cursos Presenciais da região foram alocados nos conceitos 2 ( nove cursos ) e 5 ( 13 cursos ).
Ficaram Sem Conceito ( SC ) três cursos, e nenhum recebeu o conceito 1.
Também para os cursos oferecidos na Modalidade Educação a Distância, na região Sudeste, a distribuição foi bimodal, com os conceitos 2 e 3 sendo atribuídos a nove cursos cada.
Os demais cursos receberam o conceito 4 ( oito cursos ) e o conceito 5 ( três cursos ), não tendo havido cursos Sem Conceito ( SC ) ou com o conceito 1.
As Instituições Privadas concentraram 21 dos 40 cursos participantes da região Sul, 52,5 \% do total regional.
Desses, oito cursos obtiveram conceito 3, o conceito modal.
Nessa combinação de Categoria Administrativa e Grande Região, um curso recebeu o conceito 1 ; cinco cursos receberam o conceito 2 ; dois cursos, o conceito 4 ; e ainda três, o conceito 5.
Dois cursos ficaram Sem Conceito ( SC ) nessa categoria, na região Sul.
As Instituições Públicas na região Sul participaram com 19 cursos ( 47,5 \% dos cursos da região ), e o conceito modal foi o 4, atribuído a dez cursos.
Três cursos receberam o conceito 2, quatro receberam o conceito 3, e dois, o conceito 5.
Não houve, na região Sul, curso oferecido por IES Públicas que tenha ficado Sem Conceito ( SC ) ou recebido conceito 1.
A maior parte dos cursos oferecidos na região Sul foram presenciais ( 25 dos 40 cursos participantes ).
O conceito modal para essa modalidade foi o 4, recebido por 11 cursos.
Os outros cursos presenciais, na região Sul, receberam os conceitos 2 ( um curso ), 3 ( sete cursos ) e 5 ( cinco cursos ), tendo um curso ficado Sem Conceito ( SC ).
Nenhum curso recebeu o conceito 1.
Foram 15 os cursos da Modalidade Educação a Distância participantes na região Sul, dos quais sete receberam o conceito modal 2.
Os conceitos 1 e 4 foram atribuídos a um curso cada, e cinco cursos obtiveram o conceito 3.
Um curso ficou Sem Conceito ( SC ), e nenhum recebeu o conceito 5 na Modalidade a Distância, na região Sul.
Na região Centro Oeste, oito dos 26 cursos participantes eram de Instituições Privadas ( 30,8 \% em termos regionais ).
Desses, quatro cursos receberam o conceito modal 3.
Dos demais cursos, dois receberam o conceito 2, e um, o conceito 4.
Um curso nessa categoria e região ficou Sem Conceito ( SC ), e nenhum recebeu os conceitos 1 e 5.
Quanto aos 18 cursos oferecidos por Instituições Públicas na região Centro Oeste ( 69,2 \% do total regional ), sete cursos apresentaram conceito 3, conceito modal.
Aos demais foram atribuídos os conceitos 1 ( um curso ), 2 ( cinco cursos ), 4 ( três cursos ) e 5 ( um curso ).
Um curso nessa categoria e região ficou Sem Conceito ( SC ).
No Gráfico 5.2, apresenta -se a distribuição do Conceito Enade, segundo a Categoria Administrativa da IES.
Os cursos em IES Privadas ( linha verde ) apresentam uma poligonal mais à direita do que os em IES Públicas ( linha azul ) e, consequentemente, uma distribuição de Conceito Enade com valores maiores. 5.3 CONCEITOS POR ORGANIZAÇÃO ACADÊMICA E GRANDE REGIÃO Na Tabela 5.3, encontra -se a distribuição dos conceitos atribuídos aos cursos participantes do Enade/2021, na Área de Letras-Português ( Licenciatura ), por Organização Acadêmica, segundo as Grandes Regiões brasileiras.
Dos 292 cursos de Letras-Português ( Licenciatura ) participantes, 223 eram oferecidos em Universidades, 33 em Centros Universitários, 23 em Faculdades e 13 em CEFET/IF.
Esta distribuição corresponde a, respectivamente, 76,4 \%, 11,3 \%, 7,9 \% e 4,5 \% dos cursos.
De acordo com os dados apresentados, de todos os 26 cursos avaliados com conceito 5, 21 eram vinculados a Universidades.
Os cursos vinculados a esse tipo de Organização Acadêmica tiveram o conceito 3 como conceito modal, atribuído a 80 cursos ( 35,9 \% ).
Os demais cursos vinculados a Universidades receberam os conceitos 1 ( 12 cursos ), 2 ( 64 cursos ), 4 ( 40 cursos ) e 5 ( 21 cursos, como já mencionado ).
Seis cursos ficaram Sem Conceito ( SC ).
Dos cursos em Centros Universitários, o conceito modal também foi o 3, atribuído a 14 cursos ( 42,4 \% ).
Um curso vinculado a esse tipo de Organização Acadêmica ficou Sem Conceito ( SC ), e os outros cursos receberam os conceitos 1 ( dois cursos ), 2 ( dez cursos ), 4 ( três cursos ) e 5 ( três cursos ).
Dos 23 cursos mantidos por Faculdades, nove ( 39,1 \% ) receberam conceito 2, o conceito modal.
O segundo conceito mais frequentemente atribuído aos cursos oferecidos por essa categoria de IES foi o conceito 3, recebido por cinco cursos.
Os conceitos 1 e 4 foram atribuídos a três cursos cada, e um curso recebeu o conceito 5.
Dos cursos oferecidos em Faculdades, dois ficaram Sem Conceito ( SC ).
Os CEFET/IF participaram com 13 cursos, e o conceito modal para essa organização, foi o 4, atribuído a seis cursos.
O segundo conceito mais frequente nessa Organização Acadêmica foi o 3 ( 5 cursos ).
Os conceitos 2 e 5 foram atribuídos a um curso cada, e nenhum curso ficou Sem Conceito ( SC ) ou recebeu o conceito 1.
Considerando -se, separadamente, as regiões brasileiras, verifica -se que, na região Norte, as Universidades participaram com 23 dos 27 cursos avaliados.
Nessa região, a distribuição foi bimodal, com os conceitos 2 e 3 sendo atribuídos a dez cursos cada.
Dois cursos receberam o conceito 1, e um curso, o conceito 4.
Nenhum curso ficou Sem Conceito ( SC ) ou recebeu o conceito 5 nessa combinação de categoria e região.
Os Centros Universitários da região Norte foram representados por apenas um curso, o qual obteve o conceito 2.
As Faculdades participaram com dois cursos na região Norte, os quais receberam os conceitos 2 ( um curso ) e 3 ( um curso ).
O único curso vinculado a CEFET/IF na região Norte recebeu o conceito 4.
Na região Nordeste, as Universidades participaram com 95 dos 115 cursos da Área de Letras-Português ( Licenciatura ) oferecidos.
Nenhum curso ficou Sem Conceito ( SC ).
O conceito modal 3 foi atribuído a 39 cursos.
Nove cursos obtiveram o conceito 1, e 33 cursos, o conceito 2.
O conceito 4 foi atribuído a dez cursos, e o conceito 5, a quatro cursos.
Os Centros Universitários contaram com oito cursos participantes na região Nordeste, três dos quais receberam o conceito modal 3.
Para os demais cursos, foram atribuídos os conceitos 1 ( um curso ), 2 ( dois cursos ) e 4 ( um curso ).
Um curso oferecido por Centros Universitários no Nordeste ficou Sem Conceito ( SC ) ; nenhum obteve o conceito 5.
As Faculdades foram representadas por oito cursos, na região Nordeste.
A distribuição foi bimodal, com os conceitos 1 e 2 sendo atribuídos a três cursos cada, e um curso tendo recebido o conceito 3.
Um curso ficou Sem Conceito ( SC ), e nenhum curso recebeu os conceitos 4 e 5.
Já dos quatro cursos mantidos por CEFET/IF, na região Nordeste, três receberam o conceito modal 3, e um curso obteve o conceito 2.
Na região Sudeste, as Universidades concentraram 52 dos 84 cursos de Letras-Português ( Licenciatura ) da região.
Dos cursos oferecidos em Universidades, nessa região, o conceito modal foi o 4, atribuído a 15 cursos.
Os demais cursos receberam os conceitos 2 ( dez cursos ), 3 ( 14 cursos ) e 5 ( 11 cursos ).
Dois cursos ficaram Sem Conceito ( SC ), e nenhum recebeu o conceito 1.
Os Centros Universitários participaram com 15 cursos, na região Sudeste, dos quais sete obtiveram o conceito modal 3.
O conceito 2 foi atribuído a três cursos ; o conceito 4, a dois cursos ; e o conceito 5, a três cursos.
Nenhum curso ficou Sem Conceito ( SC ) nessa combinação de organização acadêmica e região ou recebeu o conceito 1.
Já as Faculdades foram representadas por 11 cursos na região Sudeste, dos quais um ficou Sem Conceito ( SC ).
O conceito modal 2 foi atribuído a cinco cursos.
Aos demais cursos foram atribuídos os conceitos 3 ( dois cursos ), 4 ( dois cursos ) e 5 ( um curso ).
Nesse tipo de organização, nenhum curso recebeu o conceito 1.
Os CEFET/IF da região Nordeste participaram com seis cursos, e o conceito modal 4 foi atribuído a quatro cursos, tendo os demais recebido os conceitos 3 e 5 ( um curso cada ).
Nesse tipo de organização, nenhum curso ficou Sem Conceito ( SC ) ou recebeu os conceitos 1 ou 2.
Dos 40 cursos da região Sul, 33 eram mantidos por Universidades, para os quais o conceito modal foi 4, atribuído a 12 cursos.
Os demais cursos receberam os conceitos 2 ( cinco cursos ), 3 ( nove cursos ) e 5 ( cinco cursos ).
Nesse tipo de organização, dois cursos da região Sul ficaram Sem Conceito ( SC ), e nenhum recebeu o conceito 1.
Dos sete cursos participantes de Centros Universitários da região Sul, nenhum ficou Sem Conceito ( SC ) ou recebeu os conceitos 4 e 5.
Os conceitos modais 2 e 3 foram atribuídos a três cursos cada, e um curso recebeu o conceito 1.
Na região Sul, não houve cursos participantes vinculados a Faculdades e a CEFET/IF.
Na região Centro-Oeste, 20 dos 26 cursos oferecidos foram mantidos por Universidades, e, desses, oito receberam o conceito modal 3.
Para esse tipo de Organização Acadêmica, os conceitos 1 e 5 foram atribuídos a um curso cada ; o conceito 2 foi atribuído a seis cursos ; e o conceito 4, a dois cursos.
Dois cursos ficaram Sem Conceito ( SC ).
Os Centros Universitários da região Centro-Oeste contaram com dois cursos, tendo um curso sido avaliado com o conceito 2, e um curso, com o conceito 3.
Dos dois cursos oferecidos em Faculdades, na região Centro-Oeste, um recebeu o conceito 3, e o outro, o conceito 4, a mesma distribuição dos dois cursos vinculados a CEFET/IF nessa região.
No Gráfico 5.3, apresenta -se a distribuição do Conceito Enade dos cursos de Letras-Português ( Licenciatura ), segundo a Organização Acadêmica da IES.
Os cursos em CEFET/IF ( linha laranja ) apresentam uma poligonal mais à direita do que as demais, denotando uma melhor distribuição de Conceitos Enade.
No outro extremo, encontram -se os cursos mantidos por Faculdades ( linha vermelha ), apresentando as respectivas poligonais mais à esquerda e, consequentemente, distribuição de Conceitos Enade com menores valores.
Os cursos mantidos por Universidades ( linha azul ) e por Centros Universitários ( linha verde ) apresentam poligonais mais centrais e, consequentemente, a distribuição de Conceitos Enade com os valores intermediários.

% matched lemmas: elevado, enfatizar
\end{textsample}
